\chapter{Introducción}
\label{chap:introduccion}

Una red universitaria rara vez es noticia cuando funciona; lo es, casi siempre, cuando deja de hacerlo. Esta asimetría, banal en apariencia, encierra una de las paradojas centrales de la administración contemporánea de infraestructuras: aquello que sostiene la totalidad de los procesos académicos y administrativos de una institución permanece invisible hasta el instante mismo de su falla. Cabe entonces preguntarse si una universidad puede seguir confiando la continuidad de sus servicios de red a la vigilancia reactiva de sus usuarios o si, por el contrario, la complejidad alcanzada por su infraestructura exige dotarse de una capacidad propia y sistemática de observación. Tal es la interrogante que orienta el presente capítulo.

Para responderla, el capítulo procede en tres movimientos. En un primer momento, sitúa el problema en su contexto institucional y lo confronta con el estado del arte, con el fin de establecer qué se sabe y qué permanece sin resolver. En un segundo momento, formula el problema, discute las razones que justifican intervenir sobre él y delimita, a través de los objetivos y del alcance, aquello que el proyecto se compromete a realizar. Finalmente, expone los fundamentos teóricos que hacen posible la solución propuesta y que servirán de base a la metodología desarrollada en el capítulo siguiente.

\section{Contextualización}
\label{sec:contextualizacion}

La educación superior ecuatoriana ha atravesado durante la última década una transformación que no es únicamente pedagógica, sino también infraestructural. Las aulas virtuales, los sistemas de gestión académica, las plataformas de videoconferencia, los laboratorios de cómputo y los servicios administrativos en línea han dejado de ser complementos de la actividad universitaria para convertirse en su soporte cotidiano. En consecuencia, la red de datos institucional ha pasado de ser un recurso técnico entre otros a constituir la condición material de posibilidad de la docencia, la investigación y la gestión.

La Universidad Tecnológica Indoamérica (UTI), institución de educación superior particular con sede matriz en Ambato y sede en Quito, no escapa a esta tendencia. Sus procesos misionales descansan sobre una infraestructura de red de campus heterogénea, compuesta por conmutadores de núcleo, distribución y acceso, enrutadores de frontera, cortafuegos perimetrales, controladores y puntos de acceso inalámbricos, así como por servidores que prestan servicios esenciales de red, entre ellos la resolución de nombres (DNS), la asignación dinámica de direcciones (DHCP), los servicios de directorio y los portales web institucionales. Esta infraestructura, construida de forma incremental y con equipos de distintos fabricantes, es administrada por la Dirección de Tecnologías de la Información de la institución.
% Completar con datos del levantamiento (OE1): número aproximado de equipos por capa,
% fabricantes predominantes y servicios críticos inventariados.

Ahora bien, la heterogeneidad del equipamiento no constituye por sí misma una dificultad; lo es, en cambio, la ausencia de una mirada unificada sobre él. En el estado actual, la operación de la red descansa esencialmente sobre un esquema reactivo: las incidencias suelen conocerse a partir del reporte de los usuarios finales, y su diagnóstico exige consultar de forma separada las interfaces de administración de cada equipo. No se dispone de una visión consolidada y en tiempo real de los indicadores de salud de los dispositivos, del comportamiento del tráfico ni de los eventos que estos registran. A ello se suma un marco normativo exigente: la Ley Orgánica de Protección de Datos Personales \cite{lopdp2021} obliga a que cualquier tratamiento de información vinculada a los usuarios, incluida la telemetría de red, se realice bajo criterios de minimización y confidencialidad.

El contexto queda así delineado: una institución cuya actividad depende de su red, una red que crece en complejidad y una capacidad de observación que no ha evolucionado al mismo ritmo. Antes de precisar el problema que de ello se deriva, conviene examinar cómo la literatura científica y técnica ha abordado esta misma tensión.

\section{Estado del Arte}
\label{sec:estado_del_arte}

La supervisión de redes no es un campo nuevo; lo que sí es reciente es la conciencia de sus límites. Durante décadas, la administración de redes se sustentó en esquemas de sondeo periódico mediante el protocolo simple de administración de red (\textit{Simple Network Management Protocol} -- SNMP) y en pruebas básicas de alcanzabilidad mediante ICMP. Este modelo responde con eficacia a una pregunta binaria ---¿está el equipo activo o no?---, pero resulta insuficiente cuando la falla no consiste en una caída total, sino en una degradación progresiva, una saturación intermitente o una cadena de eventos cuyo origen se encuentra lejos del síntoma observado. De ahí el desplazamiento conceptual que la industria y la academia han operado desde el \textit{monitoreo} hacia la \textit{observabilidad}: ya no se trata solo de constatar estados predefinidos, sino de inferir el estado interno de un sistema a partir de sus salidas, del mismo modo en que la teoría de control definió originalmente el concepto \cite{kalman1960general, majors2022observability}.

La cuestión que guía esta revisión puede formularse así: ¿en qué medida las soluciones documentadas permiten a una institución universitaria alcanzar esa observabilidad de manera integral, segura y sostenible? Para responderla, se examinan primero los aportes relativos a cada una de las fuentes de telemetría ---flujos de tráfico, métricas de dispositivos y registros de eventos---, luego las experiencias aplicadas en entornos universitarios y, por último, el marco normativo y las plataformas de código abierto disponibles.

\subsection{Revisión de Trabajos Académicos e Industriales}

\paragraph{La telemetría de flujos: ver el tráfico sin leer su contenido.}
En el ámbito de la caracterización del tráfico, Hofstede et al. \cite{hofstede2014flow} presentaron un estudio de referencia sobre la cadena completa del monitoreo basado en flujos, desde la captura de paquetes hasta el análisis de datos, centrado en los protocolos NetFlow e IPFIX (\textit{IP Flow Information Export}, normalizado en la RFC 7011 \cite{claise2013ipfix}). Su tesis es que la agregación de paquetes en flujos, identificados por atributos comunes de cabecera como las direcciones, los puertos y el protocolo, reduce de manera drástica el volumen de datos frente a la captura completa de paquetes, lo que vuelve viable el análisis del tráfico en enlaces de alta velocidad. El argumento tiene, además, una consecuencia que trasciende lo técnico: al trabajar sobre metadatos y no sobre las cargas útiles, el análisis de flujos resulta compatible con los principios de minimización de datos personales que exige la legislación ecuatoriana \cite{lopdp2021}. Los autores no ocultan, sin embargo, la contracara de esta economía: la exactitud del análisis depende de la parametrización de los temporizadores de exportación y de las tasas de muestreo, de modo que una mala configuración puede producir una imagen distorsionada del tráfico real.

Esta limitación invita a preguntarse si los flujos bastan para comprender qué aplicaciones consumen la red. Deri et al. \cite{deri2014ndpi} responden parcialmente con nDPI, una biblioteca de código abierto para la inspección de paquetes a alta velocidad, integrada en el analizador ntopng, que identifica protocolos de capa de aplicación sin depender de números de puerto fijos. Su aporte enriquece la telemetría de flujos con una dimensión semántica; con todo, la inspección del contenido de los paquetes plantea un dilema que el presente proyecto no puede eludir, pues entra en tensión con la protección de datos personales. La lección que se extrae es, por tanto, matizada: la clasificación de aplicaciones es valiosa, pero debe limitarse a lo estrictamente necesario y apoyarse preferentemente en los metadatos de flujo.

\paragraph{Las métricas de dispositivos: el problema de la escala.}
Si los flujos describen lo que circula por la red, las métricas obtenidas mediante SNMP describen el estado de quienes la conforman. Espinel-Villalobos et al. \cite{espinel2022design} parten de una constatación precisa: en redes corporativas de gran tamaño, la recolección y el procesamiento de datos mediante SNMP pueden introducir una latencia que dificulta el monitoreo en tiempo real. Para superarla, proponen un sistema multiagente organizado en capas, en el que un agente recolector obtiene los valores de la base de información de gestión (MIB) de los equipos, un agente consolidador los procesa y los deja en un formato consumible, y un agente de aplicación los presenta como un servicio web en forma de mapa de calor. El interés de esta propuesta reside menos en su implementación concreta que en el principio que la sustenta: separar la ingesta de datos de su procesamiento y de su presentación. Su alcance, no obstante, se restringe a las métricas obtenidas por SNMP, sin integrar la telemetría de flujos ni los registros de eventos.

\paragraph{Los registros de eventos: la memoria de la red.}
Ni los flujos ni las métricas capturan, sin embargo, todo lo que ocurre en un equipo. He et al. \cite{he2016experience} evaluaron de forma sistemática distintos métodos de detección de anomalías basados en registros de sistema, tanto supervisados como no supervisados, y mostraron que los registros contienen información diagnóstica que permite detectar comportamientos anómalos, siempre que sean previamente analizados y estructurados. En la misma línea, Vaarandi y Pihelgas \cite{vaarandi2015logcluster} propusieron LogCluster, un algoritmo de agrupamiento y minería de patrones que identifica las líneas recurrentes y los eventos atípicos en registros textuales, lo que facilita al administrador la síntesis de grandes volúmenes de eventos. Ambos trabajos convergen en una premisa que resulta decisiva para el presente proyecto: el análisis de registros solo es posible si estos se conservan de manera centralizada y en un formato normalizado, como el definido por el protocolo Syslog en la RFC 5424 \cite{gerhards2009syslog}. Un registro que permanece en la memoria volátil de un conmutador desaparece, precisamente, cuando más se lo necesitaría: tras un reinicio inesperado.

\paragraph{Las experiencias en el ámbito universitario ecuatoriano.}
Las experiencias locales confirman tanto el interés como los límites de los enfoques anteriores. Obregón Maroto \cite{obregon2026implementacion}, en un trabajo de integración curricular de la Escuela Politécnica Nacional, implementó un prototipo de monitoreo de tráfico en una red local virtual, recolectando métricas mediante SNMP, almacenándolas en Prometheus y representándolas en tableros de Grafana. Su trabajo demuestra la viabilidad de construir capacidades de monitoreo con herramientas de código abierto en el contexto académico nacional. Sin embargo, su carácter de prototipo en un entorno virtual y su restricción a las métricas de SNMP dejan abiertas dos cuestiones: la integración con la telemetría de flujos y de eventos, y la validación en una infraestructura institucional real.

\paragraph{El marco normativo y las plataformas disponibles.}
Estas experiencias dispersas encuentran un marco conceptual común en la RFC 9232, \textit{Network Telemetry Framework}, publicada por el Grupo de Trabajo de Ingeniería de Internet (IETF) \cite{song2022network}. El documento organiza la telemetría según su origen en los planos de gestión, control y reenvío, y sostiene que las arquitecturas modernas deben combinar el sondeo periódico con mecanismos de envío continuo de datos. La RFC 9417 \cite{claise2023sain} complementa esta visión con la noción de aseguramiento de servicio, según la cual la salud de un servicio solo puede evaluarse a partir de un grafo de dependencias que vincule cada servicio con los subservicios y componentes de los que depende. Esta idea tiene una consecuencia práctica de primer orden: cuando un equipo central falla, las alertas de todos los equipos que dependen de él no son sino ecos de una misma causa, y una plataforma que no lo comprende satura a sus operadores de notificaciones redundantes.

En el plano de las herramientas, el ecosistema de código abierto ofrece alternativas maduras, pero no equivalentes. Prometheus se ha consolidado en entornos de microservicios por su modelo de recolección de series temporales; su integración con equipos de red, empero, depende de un exportador SNMP cuya configuración exige generar y mantener manualmente los módulos de recolección a partir de las MIB \cite{prometheus2024docs}. Zabbix, en cambio, incorpora de forma nativa el soporte de SNMPv3, el descubrimiento de bajo nivel (\textit{Low-Level Discovery} -- LLD) de interfaces y sensores, y la definición de dependencias entre disparadores, lo que permite suprimir las alertas derivadas de una falla aguas arriba \cite{zabbix2024manual}. La elección entre ambas no es, por tanto, una cuestión de preferencia, sino de adecuación al objeto supervisado.

\subsection{Comparación Sistemática de Soluciones del Estado del Arte}

La Tabla~\ref{tab:estado_arte} sintetiza los trabajos analizados, sus ventajas, sus limitaciones y el aporte que cada uno realiza a la formulación del presente proyecto.

\begin{table}[htbp]
  \centering
  \footnotesize
  \caption{Comparación de soluciones y trabajos del estado del arte en monitoreo de redes.}
  \label{tab:estado_arte}
  \renewcommand{\arraystretch}{1.1}
  \setlength{\tabcolsep}{2.5pt}
  \begin{tabular}{|p{2.1cm}|p{2.3cm}|p{2.8cm}|p{2.8cm}|p{2.8cm}|}
    \hline
    \textbf{Trabajo / Referencia} & \textbf{Enfoque y Protocolos} & \textbf{Ventajas Principales} & \textbf{Limitaciones} & \textbf{Aporte al Proyecto UTI} \\
    \hline
    Hofstede et al. (2014) \cite{hofstede2014flow} & Monitoreo de flujos NetFlow e IPFIX (RFC 7011 \cite{claise2013ipfix}). & Caracteriza el tráfico sin inspeccionar las cargas útiles; alta escalabilidad. & Sensible a temporizadores y muestreo; no supervisa la salud del hardware. & Fundamenta el módulo colector de flujos IPFIX y su compatibilidad con la LOPDP. \\
    \hline
    Deri et al. (2014) \cite{deri2014ndpi} & Inspección de paquetes con nDPI / ntopng. & Clasifica protocolos de aplicación sin depender de puertos fijos. & La inspección de contenido entra en tensión con la protección de datos. & Justifica ntopng como analizador de flujos, restringido a metadatos. \\
    \hline
    Espinel-Villalobos et al. (2022) \cite{espinel2022design} & Sistema multiagente sobre SNMP y MIB. & Separa recolección, consolidación y presentación; reduce la latencia. & Limitado a métricas SNMP; sin flujos ni registros de eventos. & Sustenta el diseño desacoplado de la ingesta y la visualización. \\
    \hline
    He et al. (2016) \cite{he2016experience}; Vaarandi y Pihelgas (2015) \cite{vaarandi2015logcluster} & Análisis y minería de registros de sistema. & Detectan anomalías y patrones no visibles en las métricas. & Requieren registros centralizados, normalizados y retenidos. & Fundamenta el concentrador Syslog conforme a la RFC 5424 \cite{gerhards2009syslog}. \\
    \hline
    Obregón Maroto (2026) \cite{obregon2026implementacion} & SNMP + Prometheus + Grafana. & Demuestra la viabilidad del código abierto en el contexto académico ecuatoriano. & Prototipo en red virtual; sin flujos ni registros de eventos. & Referente local directo que el proyecto amplía y valida en producción. \\
    \hline
    Song et al. (2022) \cite{song2022network}; Claise et al. (2023) \cite{claise2023sain} & Marcos de telemetría y aseguramiento de servicio de la IETF. & Taxonomía formal de la telemetría; grafos de dependencias. & Marcos conceptuales sin implementación de referencia. & Guía la arquitectura por capas y la supresión de alertas por dependencias. \\
    \hline
  \end{tabular}
  \par\vspace{1.5mm}
  \elaboradopor{El Investigador}
  \fuente{Revisión de literatura científica y técnica especializada}
\end{table}

\subsection{Identificación de la Brecha Tecnológica y Posicionamiento de la Propuesta}

La revisión permite extraer una conclusión que ninguno de los trabajos formula por sí solo. Cada uno ilumina una faceta de la red ---el tráfico, el estado de los equipos o sus eventos---, pero ninguno las reúne en una misma mirada. Esta fragmentación no es un defecto menor: obliga al administrador a consultar interfaces separadas y a correlacionar mentalmente datos heterogéneos, lo que prolonga tanto el tiempo medio de detección (\textit{Mean Time To Detect} -- MTTD) como el tiempo medio de reparación (\textit{Mean Time To Repair} -- MTTR).

A esta primera brecha se añade una segunda, de orden defensivo. Aunque el modelo de seguridad basado en usuario de SNMPv3 fue normalizado hace más de dos décadas \cite{blumenthal2002user}, es frecuente que los equipos de red sigan operando con SNMPv1 o SNMPv2c, cuyas cadenas de comunidad viajan en texto claro y exponen la información de topología a interceptaciones pasivas. Podría objetarse que las soluciones comerciales resuelven ambas dificultades; sin embargo, sus modelos de licenciamiento por sensor, puerto o nodo las vuelven difícilmente sostenibles para una universidad de escala media, y su carácter cerrado limita su adaptación al contexto institucional. Por último, la mayor parte de las experiencias revisadas se limita a describir el despliegue, sin someterlo a un protocolo de validación que mida cuantitativamente su capacidad de detección y notificación.

Es en la intersección de estas cuatro brechas ---fragmentación, inseguridad, costo y ausencia de validación--- donde se sitúa el presente trabajo. Su propósito no es inventar una nueva técnica de supervisión, sino integrar de manera coherente, segura y verificable aquellas que la literatura ya ha validado por separado. Queda, sin embargo, por precisar cómo se manifiestan estas brechas en el caso concreto de la Universidad Tecnológica Indoamérica.

\section{Planteamiento del Problema}
\label{sec:planteamiento}

El problema que aborda este trabajo puede enunciarse en los siguientes términos: la Universidad Tecnológica Indoamérica carece de una capacidad integrada de observación de su infraestructura y de sus servicios de red, lo que le impide detectar de forma oportuna las condiciones de falla y de degradación antes de que estas afecten a la comunidad universitaria.

Las causas de esta situación no son simples, sino que se refuerzan mutuamente. En primer lugar, la telemetría disponible se encuentra dispersa: cada equipo conserva sus propias métricas y registros, accesibles únicamente a través de su interfaz de administración, sin un repositorio común que permita su correlación. En segundo lugar, los registros de eventos de los conmutadores y enrutadores se almacenan en memorias locales de capacidad limitada, de modo que se pierden tras un reinicio o se sobrescriben con el paso del tiempo, privando al personal técnico de la evidencia necesaria para reconstruir la secuencia de una falla. En tercer lugar, no se dispone de una caracterización sistemática del tráfico que permita identificar qué aplicaciones, servicios o equipos consumen la capacidad de los enlaces. A estas causas técnicas se añade una de índole organizativa: en ausencia de umbrales definidos y de una matriz de escalamiento, la detección de incidencias depende en gran medida del reporte de los propios usuarios.

Los efectos de esta situación son igualmente identificables. El más inmediato es el aumento del tiempo que transcurre entre la aparición de una falla y su detección, durante el cual los servicios académicos y administrativos se ven interrumpidos o degradados. A ello se suma la dificultad para determinar la causa raíz de los incidentes, que conduce a diagnósticos prolongados y, en ocasiones, a soluciones que atienden el síntoma y no su origen. En el mediano plazo, la ausencia de series históricas impide planificar con fundamento la capacidad de los enlaces y la renovación del equipamiento, de modo que las decisiones de inversión se toman sin evidencia cuantitativa. Por último, la eventual persistencia de versiones no seguras de los protocolos de gestión expone información sensible sobre la topología de la red.

La pregunta que se desprende de este diagnóstico es, por consiguiente, la siguiente: ¿cómo puede la Universidad Tecnológica Indoamérica dotarse de una capacidad de observación de su red que sea integral, segura, económicamente sostenible y cuya eficacia pueda demostrarse de manera cuantitativa?

\section{Justificación}
\label{sec:justificacion}

Que un problema exista no basta para justificar una intervención; es preciso, además, mostrar que esta es pertinente, útil y realizable. La justificación del presente proyecto se construye, por ello, en tres tiempos: se expone primero la necesidad que lo sostiene, se examinan luego las objeciones que podrían oponérsele y se establece, finalmente, por qué dichas objeciones no invalidan su realización.

La primera razón que sostiene el proyecto es de orden operativo. Una plataforma de observabilidad transforma la relación de la institución con su propia infraestructura: allí donde hoy la falla se descubre a través de la queja del usuario, mañana podrá anticiparse a partir de la tendencia de un indicador. La saturación progresiva de un enlace troncal, el aumento de los errores en una interfaz, el agotamiento de la memoria de un conmutador o la indisponibilidad de un servicio de resolución de nombres son condiciones que pueden detectarse antes de que se traduzcan en una interrupción perceptible. Reducir el tiempo medio de detección equivale, en este sentido, a proteger la continuidad de la actividad académica.

La segunda razón es de orden estratégico. Las series temporales que la plataforma acumula constituyen una base de evidencia para la gestión de la capacidad, la planificación de inversiones y la negociación de los enlaces con los proveedores de servicios. La institución dejará de decidir sobre la renovación de su infraestructura a partir de percepciones para hacerlo a partir de datos. A ello se suma un beneficio en materia de seguridad y cumplimiento normativo: la adopción de SNMPv3 con autenticación y cifrado cierra un canal de exposición de información, mientras que la restricción del análisis de tráfico a los metadatos de flujo garantiza la coherencia con la Ley Orgánica de Protección de Datos Personales \cite{lopdp2021}.

Podría objetarse, sin embargo, que la institución ya dispone de las interfaces de administración de sus propios equipos y que una plataforma adicional supone un costo, una complejidad y un riesgo injustificados. La objeción merece ser tomada en serio. Es cierto que toda intervención sobre una red en producción entraña riesgos, que una plataforma de supervisión requiere recursos de cómputo y que su mantenimiento demanda competencias específicas. También es cierto que las herramientas comerciales ofrecen soluciones integradas con soporte del fabricante.

Ahora bien, ninguno de estos argumentos resiste un examen detenido. Las interfaces de administración de cada equipo, por definición, no permiten la correlación entre equipos, que es precisamente lo que el diagnóstico de fallas complejas requiere. El riesgo sobre la red en producción se acota mediante una intervención limitada a la habilitación de agentes de telemetría de solo lectura, ejecutada en ventanas de mantenimiento autorizadas y con respaldo previo de las configuraciones. En cuanto al costo, la solución se construye íntegramente con software de código abierto y sobre la infraestructura de virtualización existente, sin adquisición de licencias. Finalmente, el despliegue mediante contenedores y configuración declarativa versionada asegura que la plataforma sea reproducible y transferible al personal técnico de la institución.

La factibilidad del proyecto queda así establecida en sus dimensiones técnica, operativa y económica. Técnicamente, los protocolos que lo sustentan ---SNMPv3, IPFIX y Syslog--- son estándares abiertos ampliamente soportados por el equipamiento de red actual. Operativamente, el proyecto cuenta con el acompañamiento de la Dirección de Tecnologías de la Información de la UTI, que facilita el acceso a la red de gestión y a los recursos de virtualización. Económicamente, no requiere inversión en licencias ni en hardware nuevo.

Los beneficiarios directos del proyecto son el personal de la Dirección de Tecnologías de la Información, que dispondrá de una herramienta de diagnóstico unificada, y las autoridades institucionales, que contarán con indicadores objetivos de disponibilidad. Los beneficiarios indirectos son los estudiantes, docentes y personal administrativo de la universidad, cuya actividad depende de la continuidad de los servicios de red. Desde la perspectiva académica, el proyecto recorre el ciclo completo de la ingeniería ---análisis, diseño, implementación y validación experimental--- y aporta un referente aplicable a otras instituciones de educación superior con características similares. Justificada la intervención, corresponde ahora precisar qué se propone alcanzar.

\section{Objetivos}
\label{sec:objetivos}

\subsection{Objetivo General}
Diseñar e implementar una plataforma de monitoreo y observabilidad de la infraestructura y los servicios de red de la Universidad Tecnológica Indoamérica, basada en la recolección de métricas mediante SNMPv3, el análisis de flujos de tráfico y la centralización de registros de eventos, que permita a la institución detectar de forma proactiva las condiciones de falla y degradación del servicio.

\subsection{Objetivos Específicos}
Para alcanzar el objetivo general, se establecen los siguientes objetivos específicos, ordenados según las fases del ciclo de ingeniería (análisis, diseño, implementación y validación):

\begin{enumerate}
  \item Identificar y catalogar los activos de red y servicios críticos institucionales (conmutadores de núcleo, distribución, acceso, enrutadores de frontera, cortafuegos perimetrales, controladores Wi-Fi y servicios DNS, DHCP y Web), definiendo la matriz de métricas clave, umbrales operativos admisibles y matriz de escalamiento.
  \item Diseñar la arquitectura técnica de la plataforma bajo un modelo desacoplado y de alta disponibilidad, contemplando ingestión de SNMPv3, colector de flujos IPFIX/NetFlow, concentrador de logs Syslog conforme a la norma RFC 5424, almacenamiento optimizado para series de tiempo, capa analítica y motor de tableros.
  \item Desplegar la solución sobre un sistema operativo Linux endurecido (\textit{hardened}) mediante una arquitectura de microservicios contenerizados (Docker y Docker Compose), implementando descubrimiento dinámico de red (\textit{Low-Level Discovery} -- LLD), plantillas de supervisión específicas por fabricante/modelo y tableros interactivos orientados al servicio.
  \item Configurar el motor de alertas con umbrales multinivel (Advertencia, Promedio, Alto, Desastre), árboles de dependencia topológica para supresión de tormentas de alarmas derivadas, detección de oscilación de enlaces (\textit{flapping}), canales de notificación inmediata (Telegram y correo electrónico) y emisión automática de reportes ejecutivos periódicos de disponibilidad y consumo.
  \item Evaluar el desempeño de la plataforma en un banco de pruebas controlado mediante la inyección premeditada de eventos anómalos (caída de enlaces troncales, saturación de ancho de banda, picos de CPU, pérdida de paquetes y caída de servicios esenciales), cuantificando los tiempos de detección (\textit{Time To Detect} -- TTD) y notificación (\textit{Time To Notify} -- TTN).
\end{enumerate}

\section{Alcance del Proyecto}
\label{sec:alcance}

Delimitar un proyecto consiste tanto en afirmar lo que hará como en reconocer lo que no hará. Esta segunda dimensión no es una concesión, sino una condición de rigor: un alcance impreciso conduce a expectativas que el proyecto no puede satisfacer y a resultados que no pueden evaluarse.

En cuanto a lo que comprende, el proyecto abarca, en primer término, el levantamiento y la catalogación de los activos de red pertenecientes al perímetro acordado con la Dirección de Tecnologías de la Información: conmutadores de núcleo, distribución y acceso, enrutadores de frontera, cortafuegos perimetrales y controladores inalámbricos. Sobre estos activos se habilitará la recolección de métricas mediante SNMPv3 en modo de solo lectura, con autenticación y cifrado, a partir de los grupos estándar de la MIB-II y de las MIB propietarias de cada fabricante. Comprende, asimismo, la recolección y el análisis de flujos IPFIX o NetFlow versión 9 en los equipos de frontera, la centralización de los registros de eventos mediante Syslog conforme a la RFC 5424 y la supervisión activa de la disponibilidad de los servicios críticos de red ---ICMP, DNS, DHCP, HTTP/HTTPS y servicios de directorio---. Sobre esta base se construirán tableros diferenciados para la operación técnica y para la gestión institucional, un subsistema de alertas con supresión por dependencias y notificación por Telegram y correo electrónico, reportes periódicos automatizados y un protocolo de validación experimental mediante fallas inducidas. El proyecto concluye con la entrega de la documentación técnica y la capacitación del personal responsable.

En cuanto a lo que excluye, el proyecto no contempla la adquisición de hardware ni de licencias de software comercial, ni la modificación de la topología o de la configuración de enrutamiento de la red institucional; la intervención sobre los equipos se limita estrictamente a habilitar los mecanismos de telemetría. Tampoco forman parte de su alcance las funciones propias de un sistema de gestión de eventos de seguridad (SIEM), de prevención de intrusiones o de respuesta automatizada a incidentes, el monitoreo del desempeño interno de las aplicaciones académicas y de sus bases de datos, ni la remediación automática de fallas: la plataforma detecta, diagnostica y notifica, pero no reconfigura los equipos. Quedan igualmente fuera del alcance la supervisión de la infraestructura física del centro de datos, salvo la eventual supervisión básica de sistemas de alimentación ininterrumpida con interfaz de red, la implantación de un sistema de mesa de ayuda y la definición formal de acuerdos de nivel de servicio, cuya aprobación corresponde a las autoridades de la institución.

Así delimitado, el proyecto adquiere un contorno preciso. Resta exponer los fundamentos teóricos sobre los que descansa.

\section{Fundamentación Teórica}
\label{sec:fundamentacion}

Una plataforma de observabilidad no es una mera yuxtaposición de herramientas; es la puesta en práctica de un conjunto de conceptos, protocolos y modelos que conviene explicitar. Esta sección los presenta siguiendo el mismo recorrido que realiza un dato desde su origen hasta la decisión que motiva: se define primero el concepto de observabilidad, se examinan luego las tres fuentes de telemetría que la hacen posible y se abordan, por último, la arquitectura que las integra y los modelos de gestión que dan sentido a la información producida.

\subsection{Del Monitoreo a la Observabilidad}

El término observabilidad proviene de la teoría de control, donde Kalman lo introdujo para designar la propiedad de un sistema cuyo estado interno puede determinarse a partir de sus salidas \cite{kalman1960general}. Trasladado a la operación de sistemas informáticos, el concepto marca una diferencia sustantiva con el monitoreo tradicional. El monitoreo responde a preguntas formuladas de antemano: verifica si un valor supera un umbral o si un equipo responde a una solicitud. La observabilidad, en cambio, aspira a que el sistema produzca suficiente información como para responder preguntas que no se habían previsto, lo que resulta indispensable ante fallas nuevas o complejas \cite{majors2022observability}.

Esta distinción no implica que el monitoreo pierda su valor; implica que se convierte en una condición necesaria, pero no suficiente. La práctica de la ingeniería de confiabilidad ha propuesto, en este sentido, centrar la supervisión en cuatro señales fundamentales ---latencia, tráfico, errores y saturación--- que permiten caracterizar el estado de cualquier servicio \cite{beyer2016sre}. En una red de datos, estas señales no provienen de una única fuente, sino de tres tipos de telemetría complementarios: las métricas del estado de los dispositivos, los flujos de tráfico y los registros de eventos. Cada uno responde a una pregunta distinta y ninguno, por separado, basta para comprender el comportamiento de la red.

\subsection{Métricas de Dispositivos: el Protocolo SNMPv3}

El protocolo simple de administración de red se basa en un modelo en el que un gestor consulta a los agentes residentes en los dispositivos acerca del valor de un conjunto de variables organizadas jerárquicamente en una base de información de gestión (MIB) \cite{kurose2021computer, stallings1999snmp}. Cada variable se identifica mediante un identificador de objeto (OID) único dentro de dicha jerarquía. La MIB-II, definida en la RFC 1213, y el grupo de interfaces, actualizado en la RFC 2863, normalizan los objetos esenciales para la supervisión de red: el tiempo de actividad del equipo, el estado operativo y administrativo de cada interfaz, los contadores de octetos transmitidos y recibidos ---incluidos los contadores de 64 bits necesarios en enlaces de alta velocidad---, así como los errores y descartes \cite{mccloghrie1991mib2, mccloghrie2000ifmib}. A estos objetos estándar se añaden las MIB propietarias de cada fabricante, que exponen indicadores como la carga de procesador, la temperatura o el estado de las fuentes de alimentación.

Ahora bien, las versiones 1 y 2c del protocolo adolecen de una debilidad estructural: autentican las solicitudes mediante una cadena de comunidad que viaja sin cifrar. La versión 3 corrige esta deficiencia mediante una arquitectura modular \cite{harrington2002architecture} que incorpora un modelo de seguridad basado en usuario (USM), responsable de la autenticación y la privacidad de los mensajes \cite{blumenthal2002user}, y un modelo de control de acceso basado en vistas (VACM), que restringe las porciones de la MIB que cada usuario puede consultar \cite{wijnen2002vacm}. La autenticación se apoya en códigos de autenticación de mensajes basados en funciones resumen, extendidos a la familia SHA-2 por la RFC 7860 \cite{merkle2016hmacsha2}, mientras que la privacidad se garantiza mediante cifrado AES, incorporado al modelo por la RFC 3826 \cite{blumenthal2004aes}. La Tabla~\ref{tab:snmp_niveles} resume los niveles de seguridad resultantes.

\begin{table}[htbp]
  \centering
  \footnotesize
  \caption{Niveles de seguridad del modelo USM de SNMPv3.}
  \label{tab:snmp_niveles}
  \renewcommand{\arraystretch}{1.2}
  \begin{tabular}{|p{2.4cm}|p{2.7cm}|p{1.8cm}|p{5.5cm}|}
    \hline
    \textbf{Nivel} & \textbf{Autenticación} & \textbf{Cifrado} & \textbf{Implicación para el proyecto} \\
    \hline
    \texttt{noAuthNoPriv} & No & No & Equivalente en exposición a SNMPv1/v2c; descartado. \\
    \hline
    \texttt{authNoPriv} & Sí (HMAC) & No & Garantiza el origen de los mensajes, pero expone su contenido; insuficiente. \\
    \hline
    \texttt{authPriv} & Sí (HMAC-SHA) & Sí (AES) & Garantiza autenticidad y confidencialidad; nivel adoptado por la plataforma. \\
    \hline
  \end{tabular}
  \par\vspace{1.5mm}
  \elaboradopor{El Investigador}
  \fuente{Adaptado de las RFC 3414, RFC 3826 y RFC 7860 \cite{blumenthal2002user, blumenthal2004aes, merkle2016hmacsha2}}
\end{table}

El sondeo mediante SNMP plantea, sin embargo, un compromiso inevitable entre la resolución de los datos y la carga impuesta a los equipos. Un intervalo de consulta demasiado largo puede ocultar fluctuaciones breves; uno demasiado corto incrementa el consumo de procesador del agente y el tráfico de gestión. Este compromiso, señalado ya por Espinel-Villalobos et al. \cite{espinel2022design}, justifica la diferenciación de los intervalos de sondeo según la criticidad y la variabilidad de cada métrica.

\subsection{Flujos de Tráfico: NetFlow e IPFIX}

Las métricas de interfaz indican cuánto tráfico atraviesa un enlace, pero no qué tráfico es. Para responder a esta segunda pregunta, la telemetría de flujos agrega los paquetes que comparten un conjunto común de atributos ---típicamente las direcciones de origen y destino, los puertos de transporte y el protocolo--- en registros que resumen la comunicación \cite{hofstede2014flow}. El protocolo NetFlow versión 9, documentado en la RFC 3954, introdujo un formato basado en plantillas que permite al exportador declarar la estructura de los registros que envía \cite{claise2004netflow}. IPFIX, normalizado por la IETF en la RFC 7011, retoma este principio y lo convierte en un estándar abierto \cite{claise2013ipfix}, cuyo modelo de información, definido en la RFC 7012, especifica los elementos que pueden exportarse, como las direcciones, los puertos, el protocolo y los contadores de paquetes y octetos \cite{claise2013ipfixim}.

La exportación de flujos se rige por dos temporizadores: el temporizador activo, que fuerza la exportación periódica de los flujos de larga duración, y el temporizador inactivo, que cierra y exporta los flujos que han dejado de recibir paquetes. En enlaces de muy alto tráfico puede recurrirse, además, al muestreo de paquetes, a costa de una pérdida de exactitud. Como se señaló en el estado del arte, la calidad del análisis depende directamente de esta parametrización. Conviene subrayar, por último, que el análisis de flujos opera exclusivamente sobre metadatos de cabecera, lo que lo hace compatible con el principio de minimización de datos personales.

\subsection{Registros de Eventos: el Protocolo Syslog}

Si las métricas describen estados y los flujos describen tráfico, los registros de eventos describen acontecimientos: el cambio de estado de una interfaz, un error de hardware, un inicio de sesión administrativo o una modificación de configuración. El protocolo Syslog, en la versión normalizada por la RFC 5424, define un formato estructurado para estos mensajes que incluye una prioridad ---resultante de la combinación de una instalación (\textit{facility}) y de un nivel de severidad comprendido entre 0, emergencia, y 7, depuración---, una marca de tiempo, el nombre del equipo emisor, la aplicación y el mensaje propiamente dicho \cite{gerhards2009syslog}. Esta normalización hace posible la indexación, la búsqueda y la correlación de eventos procedentes de equipos heterogéneos.

La centralización de los registros responde a una necesidad que la literatura ha puesto de manifiesto \cite{he2016experience, vaarandi2015logcluster}: un registro que no se conserva fuera del equipo que lo genera se pierde precisamente en las circunstancias en las que resulta más valioso. Su centralización, sin embargo, exige gestionar el volumen de datos mediante políticas de filtrado por severidad y de retención temporal.

La Tabla~\ref{tab:pilares} sintetiza la complementariedad de las tres fuentes de telemetría descritas.

\begin{table}[htbp]
  \centering
  \footnotesize
  \caption{Complementariedad de las fuentes de telemetría de red.}
  \label{tab:pilares}
  \renewcommand{\arraystretch}{1.2}
  \begin{tabular}{|p{2.2cm}|p{3.3cm}|p{3.3cm}|p{3.3cm}|}
    \hline
    \textbf{Criterio} & \textbf{Métricas (SNMPv3)} & \textbf{Flujos (IPFIX)} & \textbf{Eventos (Syslog)} \\
    \hline
    Pregunta que responde & ¿En qué estado se encuentra el equipo? & ¿Qué tráfico circula y entre quiénes? & ¿Qué ocurrió y en qué momento? \\
    \hline
    Modelo de obtención & Sondeo periódico por el gestor & Envío desde el exportador & Envío desde el equipo \\
    \hline
    Puerto habitual & UDP 161 & UDP 4739 (IPFIX) o 2055 & UDP/TCP 514; TLS 6514 \\
    \hline
    Naturaleza del dato & Series temporales numéricas & Registros agregados de comunicación & Mensajes textuales con severidad \\
    \hline
    Limitación principal & No describe el contenido del tráfico & Sensible a temporizadores y muestreo & Alto volumen; requiere normalización \\
    \hline
    Normas de referencia & RFC 3411--3418, 1213, 2863 & RFC 7011, 7012, 3954 & RFC 5424 \\
    \hline
  \end{tabular}
  \par\vspace{1.5mm}
  \elaboradopor{El Investigador}
  \fuente{Elaboración propia a partir de \cite{harrington2002architecture, claise2013ipfix, gerhards2009syslog}}
\end{table}

\subsection{Arquitectura de una Plataforma de Observabilidad}

Disponer de tres fuentes de telemetría no basta; es preciso integrarlas en una arquitectura capaz de recibirlas, almacenarlas, correlacionarlas y presentarlas. Siguiendo la organización por capas propuesta por el marco de telemetría de la IETF \cite{song2022network} y el principio de desacoplamiento señalado por Espinel-Villalobos et al. \cite{espinel2022design}, una plataforma de este tipo puede describirse en cinco niveles: una capa de fuentes, constituida por los equipos supervisados; una capa de ingesta, que recibe las métricas, los flujos y los eventos; una capa de almacenamiento, optimizada para series temporales y para la indexación de registros; una capa de correlación, donde se evalúan los umbrales y las dependencias; y una capa de presentación, que materializa la información en tableros y notificaciones.

En la capa de correlación interviene un concepto decisivo: la dependencia entre componentes. Como establece la RFC 9417, la salud de un servicio depende de la de los subservicios y componentes que lo sostienen, y su evaluación exige representar estas relaciones en forma de grafo \cite{claise2023sain}. Aplicado a la supervisión de red, este principio permite que la caída de un equipo central genere una única alerta referida a la causa raíz, en lugar de una avalancha de notificaciones procedentes de los equipos que dependen de él. Zabbix implementa este mecanismo mediante dependencias entre disparadores, y lo complementa con el descubrimiento de bajo nivel, que crea automáticamente los elementos de supervisión de cada interfaz o sensor detectado \cite{zabbix2024manual}.

En la capa de despliegue, la contenerización ofrece una respuesta a dos exigencias propias de una plataforma institucional: la reproducibilidad y la mantenibilidad. Docker encapsula cada componente con sus dependencias en contenedores ligeros que comparten el núcleo del sistema anfitrión, lo que permite describir la totalidad de la plataforma en archivos de configuración declarativos susceptibles de ser versionados \cite{merkel2014docker}. De este modo, la plataforma puede reconstruirse íntegramente a partir de su descripción, lo que facilita su respaldo, su actualización y su transferencia al personal técnico.

\subsection{Modelos de Gestión: PPDIOO e ITIL 4}

La tecnología, por sí sola, no garantiza que la información producida se traduzca en decisiones oportunas; para ello es necesario inscribirla en un modelo de gestión. Dos marcos de referencia sustentan, en este sentido, el presente proyecto. El ciclo de vida de red PPDIOO, propuesto por Cisco, organiza la intervención sobre una infraestructura en seis fases ---preparación, planificación, diseño, implementación, operación y optimización--- que proporcionan una estructura metodológica coherente con el recorrido de los objetivos específicos \cite{teare2007desgn}. Por su parte, ITIL 4 ofrece las prácticas de gestión de servicios en las que la plataforma se inserta, en particular la gestión de la supervisión y de los eventos, la gestión de incidentes y la gestión de la capacidad y el rendimiento \cite{axelos2019itil}. Desde esta perspectiva, la plataforma no es un fin en sí mismo, sino el instrumento que permite clasificar los eventos según su relevancia, acelerar el diagnóstico de los incidentes y proyectar la evolución de la demanda.

Finalmente, la eficacia de la plataforma debe poder medirse. Para ello se adoptan dos indicadores que serán definidos formalmente en la metodología: el tiempo de detección (TTD), que mide el intervalo entre la ocurrencia de una falla y su detección por la plataforma, y el tiempo de notificación (TTN), que mide el intervalo entre dicha ocurrencia y la recepción de la alerta por el personal responsable. Ambos indicadores permiten transformar la afirmación de que la plataforma ``detecta proactivamente las fallas'' en una proposición verificable.

\subsection{Síntesis del Capítulo}

Al término de este recorrido, la interrogante planteada al inicio del capítulo admite una primera respuesta. La Universidad Tecnológica Indoamérica no puede seguir confiando la continuidad de sus servicios de red a la vigilancia reactiva de sus usuarios, porque la complejidad de su infraestructura excede lo que una supervisión fragmentada puede abarcar. La literatura muestra que cada una de las fuentes de telemetría necesarias ---métricas, flujos y eventos--- ha sido estudiada y validada por separado, pero que su integración en una plataforma segura, sostenible y evaluada cuantitativamente sigue siendo una tarea pendiente, en particular en el contexto universitario ecuatoriano. Los fundamentos teóricos expuestos demuestran, por su parte, que dicha integración es técnicamente posible con estándares abiertos y herramientas de código abierto.

Queda por establecer el cómo. Esa es la materia del capítulo siguiente, que describe el diagnóstico de la situación actual, la metodología adoptada y el diseño de la solución que habrá de dar respuesta al problema planteado.
